\subsection{Depth and regular sequences}

Let A be a ring, M an A-module. Recall (A, X, p. 158) that a finite sequence \((x_{1},\ldots,x_{r})\) of elements of A is said to be regular for M, or M-regular, if, for \(i=1,\ldots,r\) , the homothety of ratio \(x_{i}\) in the A-module \(\mathrm{M}/(x_{1}\mathrm{M}+\ldots+x_{i-1}\mathrm{M})\) is injective. Let \((x_{1},\ldots,x_{r})\) be an M-regular sequence; for every flat A-module N, the sequence \((x_{1},\ldots,x_{r})\) is \(M\otimes_{A}N\) regular. If \(\rho:A\to B\) is a ring homomorphism making B into a flat A-module, the sequence \((\rho(x_{1}),\ldots,\rho(x_{r}))\) is regular for the B-module \(B\otimes_{A}M\) . In particular, for every prime ideal p of A, the image in \(A_{p}\) of the sequence \((x_{1},\ldots,x_{r})\) is \(M_{p}\) -regular.

In what follows we shall consider mainly the notion of M-regular sequence in the case where the ring A is local Noetherian, the module M is finitely generated and the elements of the sequence belong to \(m_{A}\) ; the notion of M-regular sequence then coincides with that of completely secant sequence for M (A, X, p. 160, cor. 1).

PROPOSITION 7.--- Let A be a ring, J an ideal of A, M an A-module, and \((x_{1},\ldots,x_{r})\) an M-regular sequence of elements of J. One has

\[
\operatorname{prof} _ {\mathrm{A}} (\mathrm{J}; \mathrm{M}) = r + \operatorname{prof} _ {\mathrm{A}} (\mathrm{J}; \mathrm{M} / (x _ {1} \mathrm{M} + \dots + x _ {r} \mathrm{M}))
\]

and in particular \(\operatorname{prof}_{\mathrm{A}}(\mathrm{J};\mathrm{M})\geqslant r\)

The case \(r=1\) follows from remark 5 of n° 1, applied to the exact sequence

\[
0 \rightarrow \mathrm{M} \xrightarrow {(x _ {1}) _ {\mathrm{M}}} \mathrm{M} \longrightarrow \mathrm{M} / x _ {1} \mathrm{M} \rightarrow 0.
\]

The general case is deduced from it by induction on r.

THEOREM 2.--- Let A be a Noetherian ring, J an ideal of A, and M a finitely generated A-module.

\begin{enumerate}
\def\labelenumi{\alph{enumi})}
\item
  Suppose that \(\mathrm{prof}_{\mathrm{A}}(\mathrm{J};\mathrm{M})\) is finite. Then every M-regular sequence of elements of J can be extended to an M-regular sequence of length \(\mathrm{prof}_{\mathrm{A}}(\mathrm{J};\mathrm{M})\) of elements of J.
\item
  The depth of M with respect to J is the supremum of the lengths of M-regular sequences formed from elements of J.
\item
  For \(\operatorname{prof}_{\mathrm{A}}(\mathrm{J};\mathrm{M})\) to to be finite, it is necessary and sufficient that the support of M meet \(\mathrm{V}(\mathrm{J})\) , or equivalently that one has \(\mathrm{JM} \neq \mathrm{M}\) .
\end{enumerate}

Let \((x_{1},\ldots,x_{r})\) be an M-regular sequence of elements of J. One has \(r\leqslant\operatorname{prof}_{\mathrm{A}}(\mathrm{J};\mathrm{M})\) (prop. 7); suppose that the inequality is strict. Let N denote the A-module \(\mathrm{M}/(x_{1}\mathrm{M}+\ldots+x_{r}\mathrm{M})\) . One has \(\operatorname{prof}_{\mathrm{A}}(\mathrm{J};\mathrm{N})>0\) (loc. cit.), so that there exists an element x of J such that the homothety \(x_{N}\) is injective (n° 1, remark 2), that is, the sequence \((x_{1},\ldots,x_{r},x)\) is M-regular. It follows by induction that for every integer s such that \(r\leqslant s\leqslant\operatorname{prof}_{\mathrm{A}}(\mathrm{J};\mathrm{M})\) the sequence \((x_{1},\ldots,x_{r})\) can be completed to an M-regular sequence of length s, which entails assertions a) and b). Assertion c) follows from remark 1 of no. 1 and cor. 1 of th. 1 of no. 3.

COROLLARY 1. For every M-regular sequence \((x_{1},\ldots ,x_{r})\) of elements of J, the following properties are equivalent:

\begin{enumerate}
\def\labelenumi{(\roman{enumi})}
\item
  we have \(r = \mathrm{prof}_{\mathrm{A}}(\mathrm{J};\mathrm{M})\)
\item
  the sequence \((x_{1},\ldots,x_{r})\) is maximal among the M-regular sequences formed from elements of J;
\item
  the A-module \(\mathrm{M}/(x_{1}\mathrm{M}+\ldots+x_{r}\mathrm{M})\) has a nonzero element annihilated by J ;
\end{enumerate}

\[
\text {(iv)} \text {   we have   } \operatorname{Ass} \bigl (\mathrm{M} / (x _ {1} \mathrm{M} + \ldots + x _ {r} \mathrm{M}) \bigr) \cap \mathrm{V} (\mathrm{J}) \neq \emptyset .
\]

The equivalence of (i) and (ii) follows from Th. 2; the equivalence of (ii), (iii), and (iv) follows from Remark 2 of No. 1 applied to the A-module M/(x₁M + \ldots{} + xᵣM).

COROLLARY 2. --- Let A be a Noetherian local ring, and let M be a nonzero finitely generated A-module. Then

\[
\operatorname{prof} _ {\mathrm{A}} (\mathrm{M}) \leqslant \dim_ {\mathrm{A}} (\mathrm{M}) <   + \infty .
\]

Indeed every M-regular sequence of elements of \(\mathfrak{m}_{\mathrm{A}}\) is completely secant for M (A, X, p.~157, prop. 5), hence secant for M (VIII, § 3, n° 2, cor. of prop. 3); consequently its length is bounded above by the integer \(\dim_{\Lambda}(\mathrm{M})\) (loc. cit., th. 1).

